%% ==========================================================================
%%  SECTION 2 --- BACKGROUND: RECOMMENDATION AND ITS PIPELINE
%%
%%  Job: give the reader the vocabulary the taxonomy (section 3) depends on,
%%  and establish the PIPELINE framing that organises the whole survey --
%%  collection, training, serving, delivery -- because the field's approaches
%%  divide most cleanly by which stage they protect.
%%
%%  EVIDENCE: every claim here traces to notes/evidence.md.
%% ==========================================================================
\section{Background}
\label{sec:background}

\todo{Half a page of framing: recommendation is a data-mining problem whose input
      is, by construction, a record of what people consume. State that the
      survey organises the field by PIPELINE STAGE and that this section
      establishes that framing.}

\subsection{Collaborative filtering and matrix factorisation}
\label{sec:background:cf}

\todo{Neighbourhood methods versus latent-factor models. Matrix factorisation as
      the dominant latent-factor approach: given a partially observed
      user--item matrix $\ratings \in \R^{\users \times \items}$, find
      $\userembed \in \R^{\users \times \dime}$ and
      $\itemembed \in \R^{\dime \times \items}$ with
      $\ratings \approx \userembed \cdot \itemembed$.
      Cite koren2009 for the technique, hu2008 for implicit feedback.
      Note liang2018 (neural CF) and dacrema2019 (the reproducibility critique
      showing well-tuned MF remains competitive) --- the latter matters because
      it justifies why the cryptographic literature targets MF rather than
      deep models.}

\subsection{The recommendation pipeline}
\label{sec:background:pipeline}

\todo{THE ORGANISING DEVICE OF THIS SURVEY. Four stages:
      (1) COLLECTION -- ratings leave the user;
      (2) TRAINING -- a model is fitted to the collected data;
      (3) SERVING -- scores are computed for a user and a ranking produced;
      (4) DELIVERY -- the user actually fetches the recommended item.
      Emphasise that these are separable, that different papers protect
      different subsets, and that a system is only as private as its weakest
      stage.}

\begin{figure}[t]
\centering
\todo{FIGURE (TikZ): the four-stage pipeline as a horizontal flow, with what an
      adversary observes annotated beneath each stage. This figure carries the
      survey's organising idea and is referenced from sections 3, 5, 6 and 7.
      Keep it schematic --- boxes, arrows, and leakage labels.}
\caption{The recommendation pipeline and what is exposed at each stage.}
\label{fig:pipeline}
\end{figure}

\subsection{What the adversary sees}
\label{sec:background:adversary}

\todo{Stage by stage, and this is where the motivation is earned rather than
      asserted:
      COLLECTION -- raw ratings.
      TRAINING -- the data set, and intermediate model state.
      SERVING -- which user asked, and what was scored highly.
      DELIVERY -- which item was fetched.
      Then the empirical evidence that this is not hypothetical:
      narayanan2008 de-anonymised the Netflix Prize data set against IMDb;
      calandrino2011 showed that the PUBLICLY VISIBLE OUTPUTS of a recommender
      --- the ``customers who bought this also bought'' lists --- leak
      individual transactions. calandrino2011 is the strongest single
      motivation in this survey: it attacks the system's intended output, not
      a leak or a breach.}

\subsection{Utility, and why it constrains the cryptography}
\label{sec:background:utility}

\todo{Brief. Recommendation quality is measured by ranking metrics (nDCG@$k$,
      Recall@$k$); a private recommender that recommends badly has solved
      nothing. This constrains protocol design --- e.g. fixed-point arithmetic
      must not degrade the ranking --- and it is why section 8 insists that
      quality be reported alongside cost.}
